\quellseite{557}
\chapter[§ 36. TEACHING THE REMAINING MATERIAL]{§ 36. METHOD OF TEACHING THE REMAINING MATERIAL}

Compare the literature in chapters 26, 29, and 30.

The \emph{hymn} cannot generally be accorded independent treatment while our school system remains subject to present conditions. We shall have to be content with illuminating it, as occasion may require, now as a whole now in part, in connection with instruction in Biblical History, the Catechism, and Church History. When it is there brought in at the right moment, read by the teacher in an impressive manner, and the pupils are then drilled in it in order to read it well and with expression, and, moreover, are given a short explanation of those words that need an explanation, all is virtually done that is necessary for its comprehension (p. 379). When a festival approaches, selection is made of some striking festive hymn; festive feelings are aroused in the children's minds; they are shown how those festive feelings are given expression in the hymn, everything unintelligible meanwhile being explained as briefly as possible; it is then read in an impressive manner, at first as a whole, then the first stanza; the class is thereupon taught to read the first stanza with due attention to expression; last of all, the teacher sings or plays the melody and then practices it with the class. By this time the children almost remember the words, so that it will soon be possible, perhaps in the next lesson already, to have a few stanzas sung from memory. Frequent reviews, especially through singing, will make the children's knowledge of it secure. That no hymns should be assigned for memorizing without having been treated in the manner here set forth, should appear self-evident. Should there be time (cf. also p. 431) to devote more attention to
\quellseite{558}
a hymn, let care be taken to avoid the superabundance of other material under which teachers are wont to bury it. Even where more time is available, it will be sufficient to treat the hymn according to Schueren’s suggestions (pp. 119 ff.), whom we quote in the Practical Examples given at the end of this section.

In what manner the reading of the Bible is to be connected with the instruction in Biblical History and in Catechism, has already been shown (pp. 371 f., 416, 532 f.). Where Bible reading is treated as an independent branch, the following order will be found to answer the purpose best: 1. The teacher, using the acroamatic or erotematic method, establishes the context to the extent that is desirable and practicable, with frequent recapitulations covering the whole Gospel or Epistle; 2. when such review does not suffice, he prepares the ground for the new section by pointing out the fundamental thoughts and reading the whole with due expression; 3. he lets the children read the section (those rather intelligent are called upon first, then the backward ones; the more important parts are read in unison); he intersperses the reading with the necessary explanations (only such), by paraphrasing difficult expressions, dissolving involved sentences, explaining geographical, archeological features, etc.; 4. he has the class trace the fundamental thought, connects it with something already known, etc., and concludes with the question: “What does this section teach you, children?” Especially when the Epistles are read, one must guard against losing oneself in details, but try to discover and pursue the thread of the discourse: when the Gospels are read, it is advisable occasionally to read the parallel sections, or to have the brief account in the section read supplemented by additional features gleaned elsewhere in Biblical History, and thus to explain Scripture by itself (cf. Schueren, p. 47 ff.).

\quellseite{559}
In regard to \emph{Church History} we generally have to confine ourselves to a rousing narration of epochal events, with special emphasis upon educative features, and to drill-questions. Where the requisite amount of time is available, one may readily accommodate himself to the “formal steps”, among which presentation and association require the greater amount of care.—In the \emph{Young People’s Society} more freedom is naturally vouchsafed. While the acroamatic form of instruction is bound to preponderate here, the need for self-activity should not be overlooked. Where a City Library or a Young People’s Library renders the requisite literature available, historic sections, delineations of characters, etc., may be read, followed by brief essays on the subject, which, in turn, can be made the starting-point of discussion. This applies especially to such questions regarding church and religion which are at the time being of general interest (p. 432). Also questions relating to church usage and discipline (p. 433), difficult passages of Scripture, etc., can be illuminated in the manner suggested. It falls to the leader to remain in control of the discussion and, at the end, to give a lucid, concise, and impressive generalization. But the converse rule, too, will be found available: the leader delivers a lecture, and the discussion follows. When, for instance, in connection with the biography of Athanasius, the Nicene Creed becomes the topic for discussion, the young people may be invited to show at the next meeting in how far inalienable elements of the Christian faith were at stake; in how far the testimony of Scripture favored Athanasius; in what churches views akin to those of Arius are found today. Or, the origin of the Papacy is pictured along purely historic lines, without anticipating a verdict, whereupon attention is directed to such Scripture passages as Matt. 16 and others and the proof furnished why these can in no wise support the papal idea. Or, the origin of monasticism has been made the subject of
\quellseite{560}
a lecture; this is followed by a discussion in which the fundamental traits of a truly Christian ideal of life, with its stress upon the duties of the secular calling, are substituted for the monastic ideal. The fact that, under Constantine, the Church becomes a State Church, supplies an opportunity to exchange opinions upon the relation between State and Church. A description of the ruin of the Medieval Church may issue in a series of questions regarding the scripturalness of the distinctive doctrines of the Catholic Church or of the legitimacy of the charge made against the Reformation that moral life did not decline until after its advent. A study of the life of Luther may give rise to a discussion of the fundamental difference between the Lutheran and the Roman Church, of the meaning and the bearings of the doctrine of justification by faith alone, of the relation of Luther to Paul, of the difference in the conception of faith in the Lutheran and the Roman camp respectively (cp. especially \emph{Reu, Luther’s Life, Sketched for Young People’s Societies, and the Necessary Directions for General Discussion Appended}, 1917). A lecture on the Thirty-years’ War will furnish an opportunity to discuss the ultimate aims of Rome in regard to our country, etc.

\quellseite{561}
\chapter{§ 37. METHOD AND PERIOD}

F. Zange, \emph{Schulagende}, 1893. O. Zuck, \emph{Schulandachten fuer das ganze Jahr}, 1861. Knaudt, \emph{Schulandachten}, 1910.

There must be careful preparation for every catechetical period. That the catechist is constantly occupied with the study of catechetical material and will not ignore new facts as to method (cf. ch. 38), is a self-evident postulate; but each individual period, too, requires preparation. The catechist must have clearness regarding the ground to be covered, not only in general but down to the last detail; regarding the aim to be fixed, both the outer and the inner aim; and the way to reach it. While for the beginner, it is advisable, under certain circumstances necessary, to write down his catechesis in minutest detail, it will later be sufficient to draft a sketch; eventually a careful mental survey of the subject matter will prove adequate. The catechist will do well by providing a notebook both for his instruction in Biblical History and the Catechism, in which should be entered whatever he has found valuable for the several steps of preparation, statement of aim, presentation, penetration, application (Biblical History) and, likewise, for the intuitional basis, comparison, valuation, application (Catechism). Entering the class-room thus equipped, he will also recall what the next paragraph has to say about his personal relation to the children and the Word he is to offer them; then he will invoke the blessing of God and address himself to his task.

Every religious period is to be opened with prayer. Whether the class speaks in unison with the catechist or the catechist speaks the prayer alone, it should always be a real prayer, brought before God with a concentration of all the energies of the soul. The example of the teacher is of
\quellseite{562}
great importance here, the more so as many children are under no constraint at home in this respect, often even without opportunity of hearing prayer. Prayer, therefore, is not the least part of the education to be imparted. It is, furthermore, often advisable to connect a hymn with the opening prayer, which should be sung with vim and vigor; while it may bear in point of content upon the character of the lesson, this is by no means essential. Often one may simply sing a few stanzas of a morning hymn and read the morning prayer contained in the Catechism. Our fathers knew why they instituted this practise; often, when the lesson took place in the afternoon, it was brought to a close with the evening prayer.

In the partial parish or catechumen school, the opening prayer will at once be followed by Biblical History, when the new story should be presented and afterward read by the class several times from the textbook; an opportunity will thus be presented for a brief explanation of an occasional obscure word. After a Penmanship period or the like, arranged for the purpose of resting the children’s minds, there follows the penetration and application of the story which was presented on the day before, and meanwhile has been studied by the children at home. A period as provided by the curriculum of the Grammar school, and the singing of a hymn, brings the forenoon to a close. The afternoon is entirely devoted to branches of study prescribed by the curriculum of the public school. On the next day instruction in the Catechism takes the place of Biblical History. The opening with song and prayer is followed by a recitation in unison of the whole Chief Part of which a part is to be treated in detail on that day. In this the teacher will insist upon scrupulously exact, precisely articulated, but at the same time euphonious, speaking. Then comes the drill in the lesson of the previous Catechism period, to be followed
\quellseite{563}
at once by the new catechization. The rehearsal of the material to be memorized is connected with penetration and application (Biblical History), or with the drill (Catechism). The preparation for the new material to be memorized is connected with presentation (Biblical History), or with the new catechization (Catechism). Should one say, This arrangement does not allow a successful covering of the secular studies prescribed by the regular curriculum of the public school, we answer from our own experience and that of many others that this objection is not valid. In rural districts parents often are satisfied if their children finish Grammar school a year or so later. In this case instruction in Biblical History and Catechism can be given on each school-day—the one in the forenoon, the other in the afternoon, and the remaining hours can successfully be devoted to written work based upon the religious lesson, to Bible reading, Bible Geography, Church History, hymn drill and explanation.

The curriculum for the \emph{Saturday and summer school} has already been determined (p. 429 f.). In the \emph{regular parish school} instruction will follow the same general lines as in the partial parish or catechumen school. For the \emph{Sunday school} the following order is recommended: 1. Liturgical opening; 2. the leader tells the whole Sunday school or, at least, the whole Second Course (children from 7–12 years of age), the new story as clearly as possible, brings in at the proper place the material to be memorized, and lets the presentation distinctly issue in the golden text or the principal hymn stanza. 3. Thereupon the pupils retire with their teachers into their respective class-rooms, where the story that had been presented by the leader is read from the Biblical History and the material to be committed to memory repeatedly rehearsed (speaking in unison!); thereupon the teachers review the lesson of the previous Sunday by means of the review questions appended to the lesson. 4. Finally,
\quellseite{564}
the whole school with all its teachers assembles again, and the service is concluded according to the liturgy.—The \emph{Young People’s Society} has an order of services as fixed by its constitution: a hymn, prayer, the Bible lesson with informal discussion, hymn, lecture with discussion, closing hymn—these are the fundamental features generally observed.

\quellseite{565}
\chapter{§ 38. METHOD AND MAN}

Augustine, \emph{De catechizandis rudibus}, comp. ch. 7. G. v. Zezschwitz ii 2, 1, § 30 and 31, 1869. Th. Harnack, \emph{Katechetik}, pp. 114–118. F. Zange, pp. 46–57. A. Eckert, pp. 175–195. J. Berndt, pp. 73–78. E. Chr. Achelis, pp. 345–356.

But what kind of man should the catechist be, if he shall succeed in imbedding and anchoring the catechetical material in the intellectual life of the catechumen, in arousing his interest, and in guiding him in the paths in which the Spirit, in due time, shall lead him to personal faith? While much of what has been said in the second part of this book might be resumed here, some features of superlative importance remain to be pointed out. What has been said in regard to method finds its supplement in what is here said about the person of the catechist.

The first requirement of the catechist is a good \emph{education}. Theological education first of all. That he knows better and more accurately than the young what things are to be taught, is not sufficient; there must be mastery of the \emph{material}; that is, it should be known in its organic connection and be a real mental possession, so that all details may be given in their central bearings. Only one who has a systematic comprehension of Christian doctrine is able to recognize the relative unimportance or importance of any doctrine; only he is conversant with the threads that connect the individual doctrine with the sum of truth; nothing is viewed by him singly, but always as a part of the whole. Where systematic training is absent the catechist will lose himself in the mass of detail; will not know how to ask questions properly; nor shall he be able to rate the answers according to their content. He stumbles from one embarrassment into the other, and impresses one as a man who has
\quellseite{566}
no sure ground under his feet. He who would be an efficient catechist should accordingly study dogmatics. While it is perfectly true that catechetical instruction should never bear a dogmatical complexion, it is just as true that the safe foundation required for instruction is an outgrowth of a diligent study of dogmatics: what the spine is to the human body, dogmatics is to catechetical instruction. But inasmuch as all his teaching material should be based upon Scripture, and catechetical material has been taken from Scripture, the continued study of dogmatics should be accompanied by a pursuit of exegesis. Only as he goes down, again and again, into the ocean of redemptive truth gathered in the Bible, does his instruction possess the required vigor and, with all freedom, the stamp of noble dependence. Alongside of theological education dialectic versatility is required. The development and the drill particularly are dialectic exercises. But the dialectic process has to be learned and practiced. The laws that control thinking have to be understood if one wishes to think logically himself and teach others to do likewise; the mind has to acquire the necessary agility and alertness. Attentive hearing of model catecheses, and, wherever that is possible, diligent study of good printed catecheses is, in addition to the study of logic, the best aid in the acquisition of dialectic skill. The beginner had best put down his catecheses in writing; even the mature catechist should not shrink from mentally fixing the outline of the catechization—according to its major and minor points. While this is desirable from the viewpoint of dialectic clearness and skill, it is even more so from that of the importance and responsibility of such instruction. Finally, the main results of the psychological science should not be foreign to the catechist; by constantly taking into consideration the psychic life of the children he will make his task easier for himself and more successful.

\quellseite{567}
The second requirement to be met by the catechist is \emph{love to the children}. The heart of the child, and that of the adult catechumen likewise, is hungry for love; and love is what the catechist owes his pupil as a Christian, still more as a trainer. Without love, the training effort is without a soul. To be sure, the love here spoken of cannot be that earthly, effeminate, spineless love which rates the children according to their several degrees of endowment and rank, becomes extremely partial in consequence, and, for this reason, arouses in the pupils envy, jealousy, and a feeling of being slighted. The love of which we speak is rather that which is brought to bear upon all alike and which has its ground in the thought that they are children of God that have been entrusted to him. Such love is not blind to the sins and evil propensities of a child; it rather clears one’s vision to see them and generates the power to fight them with irrepressible determination. Such love can inflict punishment—painful punishment; but the punishment is inflicted for the purpose of correction, and the child feels that love has inspired it. Though the severest reproof be administered, and even corporal punishment be resorted to—a penalty never to be inflicted except in extreme cases, the child searches for love in the eye of its monitor nevertheless; and when it is beheld, together with sadness over the necessity of administering punishment, he will not only willingly submit to it but permanently benefit by it. Such love finds the tone that reaches the heart; from it flows that \emph{hilaritas} already demanded by Augustine, the counterpart of all acerbity and melancholy—a ray of vernal sunshine in the school-room. Such love will never weary of descending to the intellectual status of the child and becoming a child among children. From such love flows that faithfulness which does not cease improving one’s methods and raising the instruction to a higher level of efficiency; which does not despise
\quellseite{568}
the task of teaching children; which works with gentleness and humility in the field of catechization, though the opportunity to shine before men is utterly absent and one’s efforts are seldom controlled by any one, knowing that the Church is blessed by one’s labors; which does not permit itself to be embittered or to be turned aside from its care for the young by thoughtless men; which always bears in mind the future of the Church and the welfare of the children, and conscientiously “redeems” the time of youth, when the spirit can be molded like wax (Eph. 5:16). Such love will not forget prayer for the catechumens, least of all for the frivolous and obstinate.

The third requirement is a \emph{believing personality, sanctified by the Holy Spirit}. If the souls of the children are to yield a real anchorage to those sacred truths; if the will also is to be moved by them, the influence brought to bear upon the intellect through a clear and intelligible presentation and a lucid and logically correct development must be supplemented by a stir of the feelings; for the way to the will is found only through the intellect and the feelings. Not only interest in the truth, but also the sympathetic, ethical, and religious interests, for which the feelings serve as bearer, require cultivation. This, to be sure, is already in a measure accomplished when the catechist, during the process of instruction, brings out how one ought to share in the joy and the grief of his fellow-men, how goodness invariably finds its reward, and wickedness its punishment; when he makes the souls of the catechumens acquainted with the characters of the Bible and of Church History—men and women morally holy, meek in the presence of God, upheld by faith; but most powerfully such feelings and interests are aroused, and most enduringly the will is prompted to similar conduct when all this is surmounted by the example of the teacher as of a man intent upon the truth alone, \emph{full}
\quellseite{569}
\emph{of sympathy for whatever of grief or glee befalls his fellow-men, a champion of the right and quickened by sincere godliness}. There is a movement with which the world of pedagogy is now astir, which advocates the axiom: Religion is taught, not by word, but by example. That is a dangerously one-sided truism, an undervaluation of the Word and of knowledge, a slighting of the historic side of the Christian religion: but the element of truth in this movement, that, in the achievement of enduring effects, felt also by the will, the living example of the teacher whose life is an embodiment of Christianity, is necessary—this is a fact hardly capable of exaggeration. All strictness of discipline, all didactic accuracy and punctiliousness, all dialectic skill and faithfulness in preparation, will prove utterly inadequate where such personal influence is lacking; where the children are not made to feel that what the teacher offers to them appears precious, holy, important also to him—is, in truth, the greatest treasure and the highest norm of his life. The instruction imparted by such a catechist is never dry: it is ever fresh and vivid; for everything said by him is constantly born afresh, as it were, from his life of faith; a grasping of truth constantly renewed; a subjection to the Word willing and incessant. They are not empty formulae which he handles like lifeless tools; facts and verities they are to which the soul draws near for sustenance! Many a thing once upon a time ushered into the soul of the child through careful development may be forgotten when the graciously earnest and genuinely godly form of the old teacher and shepherd still abides in the memory as a vital force. Such a teacher cannot fail to secure and maintain that degree of authority over the children which is his due. Little need will he have for pains and penalties to enforce authority: his very presence and attitude will fascinate the obedient children and impel the obedience of the others. Without exception they will all feel: Here
\quellseite{570}
stands one who has the right to teach us; for he does not teach what he himself fails to believe and to practice: a Johannine personality is he; to lead us to Christ is his aim.

Finally, the catechist as well as the homilist should have a \emph{living faith in the creative power of the Word of God}, which is the merchandise and tool of his trade (Matt. 13: 33; Rom. 1: 16; 2 Cor. 2: 15–17). Likewise he should possess that patience and perseverance which can wait until the Word of God imbedded in the heart begins to stir and reveals its vital power, saving or judging (Mark 4: 26–29). It is this thought of the divine power in the Gospel; of the receptivity of the chatechumens in the time of their youth; of Baptism, through which those children objectively became children of God; and of the transcendent importance of his labors for the future of the Church, which, again and again, fills the catechist with strength and cheer for his calling. And the same thought tides him over the failure which, upon more occasions than one, confronts him.

This is the ideal. A zealous suppliant at the heavenly Father’s feet shall reach it more and more. He shall be enabled rightly to anchor in the whole inner life of the young those holy truths which sustain the life of the mature congregation, to arouse an interest therein and desire for participation in the divine services, so that, when the Holy Spirit, in His own time and hour, arouses them to faith by means of those sacred truths or of a renewed presentation of the Word, they shall be able to exercise their rights and duties as members of the mature congregation.
